\section*{Tóm tắt đồ án}
\addcontentsline{toc}{section}{Tóm tắt đồ án}
Trong kỷ nguyên chuyển đổi số, dữ liệu đã trở thành tài sản cốt lõi quyết định hiệu quả quản lý của các cơ sở giáo dục đại học. Tuy nhiên, thực trạng phân mảnh thông tin giữa các phòng ban do sử dụng nhiều phần mềm độc lập đang gây ra những khó khăn lớn trong việc đồng bộ, dẫn đến sự trùng lặp và không nhất quán dữ liệu.
Để giải quyết bài toán này, đồ án tập trung nghiên cứu và hiện thực hệ thống Trục tích hợp dữ liệu trung tâm dành cho trường đại học. Hệ thống được thiết kế theo kiến trúc trục nan hoa (Hub-and-Spoke), xây dựng trên nền tảng mã nguồn phân mô-đun của NestJS kết hợp với nguồn Kong API Gateway để kiểm soát lưu lượng và đảm bảo an ninh. Trục tích hợp đóng vai trò làm trung gian tiếp nhận dữ liệu từ các API nguồn, đối soát cấu trúc dựa trên đặc tả siêu dữ liệu tập trung (Schema Registry), tiến hành làm sạch, chuẩn hóa và ghi nhận vào cơ sở dữ liệu dùng chung PostgreSQL.
Trong quá trình xây dựng, nhóm thực hiện đã tập trung giải quyết các bài toán thực tế phát sinh trong tích hợp dữ liệu quy mô lớn thông qua việc thiết lập hệ thống phòng thủ xung đột dữ liệu hai lớp, tự động phát hiện biến động cấu trúc nguồn (Schema Drift), xử lý các bản ghi mồ côi và triển khai cơ chế sao lưu hai tầng linh hoạt sử dụng MinIO nội bộ kết hợp lưu trữ đám mây AWS S3. Bên cạnh đó, hệ thống tích hợp giải pháp phân quyền truy cập linh hoạt theo vai trò kết hợp lọc phạm vi bảng bằng biểu thức chính quy.
Kết quả thực nghiệm cho thấy hệ thống vận hành ổn định và đạt hiệu năng tối ưu. Chiến lược đồng bộ tăng tiến giúp giảm thiểu đến 96\% tải lượng truyền tải mạng và thao tác ghi cơ sở dữ liệu so với phương pháp đồng bộ toàn phần. Toàn bộ quy trình kiểm thử tự động và đường ống CI/CD đã được cấu hình hoàn chỉnh thông qua GitHub Actions, cho phép đóng gói container Docker và triển khai tự động, an toàn lên máy chủ AWS EC2.

\newpage

\section*{Abstract}
\addcontentsline{toc}{section}{Abstract}
In the era of digital transformation, data has emerged as a core asset determining the management efficiency of higher education institutions. However, the fragmentation of information across departments due to the use of independent software systems poses significant synchronization challenges, leading to data redundancy and inconsistency.
To address this issue, this capstone project researches and implements a centralized Data Integration Hub tailored for universities. The system is designed based on a Hub-and-Spoke architecture, built upon a modular monolith backend using NestJS, and integrated with Kong API Gateway for traffic management and security. The integration hub acts as a mediator that consumes data from source APIs, validates data structures against a centralized Schema Registry, and performs cleaning and normalization before securely synchronizing the data into a shared PostgreSQL master database.
During development, the project focused on resolving critical real-world challenges in large-scale data integration. This was achieved by establishing a two-layer data collision defense system, automatically detecting and handling structural changes in source systems (Schema Drift), managing orphan records, and implementing a flexible two-tier backup mechanism using local MinIO and AWS S3 cloud storage. Additionally, the system incorporates a granular, role-based access control solution combined with table-scope filtering utilizing regular expressions.
Experimental results demonstrate that the system operates stably with optimal performance. The implemented Incremental Sync strategy reduces network bandwidth and database write operations by up to 96\% compared to full synchronization. Finally, the automated testing workflow and CI/CD pipelines have been fully configured via GitHub Actions, enabling Docker containerization and secure, automated deployment to AWS EC2 instances.